%% ======================================================================
%% APPENDIX G — 10-PHASE RESEARCH ANALYSIS FRAMEWORK
%% Applied to All Seven Disease Domains
%% ======================================================================

\chapter*{Appendix G: Ten-Phase Research Analysis Framework}
%% TOC entry is in main.tex (Appendix G label)
\label{app:analysis_framework}

\normalsize\normalfont\color{black}

This appendix documents the complete ten-phase research analysis framework applied to all seven disease domains (Schizophrenia, Epilepsy, Stress, Autism Spectrum Disorder, Parkinson's Disease, Depression, and Cancer). Each phase is presented with its master configuration, per-disease results, decision points, and quality checks. All figures use TikZ; all tables follow the thesis formatting standard (booktabs, full-width, \texttt{ggunavy} header accent).

%% ----------------------------------------------------------------------
%% G.0  FRAMEWORK OVERVIEW
%% ----------------------------------------------------------------------
\section{Framework Overview and Phase Flow}
\label{sec:framework_overview}

% MOVED TO CHAPTER 3 MAIN TEXT (Mar 2026)
\iffalse
\needspace{20\baselineskip}
\begin{figure}[!htbp]
\centering
\begin{tikzpicture}[
    phase/.style={rectangle, rounded corners=4pt, draw=ggunavy, fill=ggunavy!8, minimum width=3.2cm, minimum height=1.1cm, font=\small\bfseries, text=ggunavy, align=center},
    arr/.style={-{Stealth[length=4mm, width=3mm]}, very thick, ggunavy},
    node distance=0.8cm
]
%% Row 1: Phases 0-4
\node[phase, fill=blue!10] (p0) at (0,0) {Phase 0\\[-1pt]{\small Research Design}};
\node[phase, fill=cyan!10] (p1) at (3.8,0) {Phase 1\\[-1pt]{\small Data Preparation}};
\node[phase, fill=green!10] (p2) at (7.6,0) {Phase 2\\[-1pt]{\small Descriptive Analysis}};
\node[phase, fill=yellow!12] (p3) at (11.4,0) {Phase 3\\[-1pt]{\small Reliability \& Validity}};

%% Row 2: Phases 4-7
\node[phase, fill=orange!12] (p4) at (0,-2.2) {Phase 4\\[-1pt]{\small Factor Analysis}};
\node[phase, fill=pink!15] (p5) at (3.8,-2.2) {Phase 5\\[-1pt]{\small Regression + SEM}};
\node[phase, fill=violet!10] (p6) at (7.6,-2.2) {Phase 6\\[-1pt]{\small Mediation + Mod.}};
\node[phase, fill=teal!12] (p7) at (11.4,-2.2) {Phase 7\\[-1pt]{\small QCA + TISM}};

%% Row 3: Phases 8-10
\node[phase, fill=lime!12] (p8) at (1.9,-4.4) {Phase 8\\[-1pt]{\small Qualitative Analysis}};
\node[phase, fill=brown!10] (p9) at (5.7,-4.4) {Phase 9\\[-1pt]{\small Integration}};
\node[phase, fill=red!10] (p10) at (9.5,-4.4) {Phase 10\\[-1pt]{\small Conclusion}};

%% Arrows
\draw[arr] (p0) -- (p1);
\draw[arr] (p1) -- (p2);
\draw[arr] (p2) -- (p3);
\draw[arr] (p3.south) -- ++(0,-0.4) -| (p4.north);
\draw[arr] (p4) -- (p5);
\draw[arr] (p5) -- (p6);
\draw[arr] (p6) -- (p7);
\draw[arr] (p7.south) -- ++(0,-0.4) -| (p8.north);
\draw[arr] (p8) -- (p9);
\draw[arr] (p9) -- (p10);

%% Feedback loop
\draw[arr, dashed, ggunavy!50] (p10.south) -- ++(0,-0.6) -| node[below, font=\small, pos=0.25, text=ggunavy!60] {Iterate if needed} (p0.south);

\end{tikzpicture}
\caption[Ten-phase research analysis framework]{Ten-phase research analysis framework: end-to-end flow from research design through data preparation, descriptive analysis, reliability testing, factor analysis, SEM, mediation/moderation, alternative models (QCA/TISM), qualitative analysis, integration, and final conclusion. A feedback loop supports iterative refinement.}
\label{fig:10phase_flow_app}
\end{figure}

\needspace{8\baselineskip}
\noindent Table~\ref{tab:app_10phase_master} summarises ten-phase framework master configuration --- objectives, methods, and outputs per phase for appendix review.

{\tablestripes\tablefontsize
\begin{xltabular}{\textwidth}{@{} c l L L L @{}}
\caption[Ten-Phase Framework Master Configuration]{Ten-Phase Framework Master Configuration --- Objectives, Methods, and Outputs per Phase}\label{tab:app_10phase_master} \\
\toprule
\thead{Phase} & \thead{Name} & \thead{Objective} & \thead{Key Methods} & \thead{Output} \\
\midrule
\endfirsthead
\multicolumn{5}{@{}l}{\textit{(\,Tab.~\ref{tab:app_10phase_master} continued from previous page\,)}} \\
\toprule
\thead{Phase} & \thead{Name} & \thead{Objective} & \thead{Key Methods} & \thead{Output} \\
\midrule
\endhead
\bottomrule
\endlastfoot
\tablestripes
0 & Research Design & Define hypotheses, variables, and framework & Literature review, conceptual modelling & Hypotheses H1--H5, research model \\
1 & Data Preparation & Clean, screen, and validate datasets & Missing value analysis, outlier detection, normality & Clean analysis-ready datasets \\
2 & Descriptive Analysis & Summarise sample and variable behaviour & Frequencies, means, SD, distributions & Sample profile, construct summaries \\
3 & Reliability \& Validity & Validate measurement quality & Cronbach's $\alpha$, CR, AVE, HTMT & Reliable and valid constructs \\
4 & Factor Analysis & Identify and confirm latent structures & EFA (Promax), CFA (model fit) & Validated measurement model \\
5 & Regression + SEM & Test hypothesised relationships & Multiple regression, PLS-SEM & Path coefficients, R$^2$, hypothesis results \\
6 & Mediation + Moderation & Test indirect and conditional effects & Bootstrapping, interaction terms & Indirect effects, conditional effects \\
7 & QCA + TISM & Validate alternative causal paths & fsQCA truth tables, TISM hierarchy & Multiple paths, structural hierarchy \\
8 & Qualitative Analysis & Explain ``why'' behind numbers & Thematic coding, discourse analysis & Themes, narratives, triangulation \\
9 & Integration & Synthesise all findings & Triangulation, model building & Final unified framework \\
10 & Conclusion & Summarise and recommend & Objective validation, contribution mapping & Contributions, limitations, future work \\
\end{xltabular}}
\fi

% [SPACE-FIX] clearpage removed to eliminate empty page

%% ======================================================================
%% G.1  PHASE 1 — DATA PREPARATION (Per Disease)
%% ======================================================================
\section{Phase 1 --- Data Preparation}
\label{sec:phase1_data_prep}

Phase 1 ensures each disease dataset is clean, valid, and analysis-ready. This section documents the data preparation steps applied to all seven domains.

\subsection{Phase 1 Master Configuration}

Table~\ref{tab:phase1_config} documents the master configuration for Phase~1, specifying the objectives, inputs, outputs, tools, and quality gates that govern data preparation across all seven disease domains.
\needspace{8\baselineskip}
{\tablestripes\tablefontsize
\begin{xltabular}{\textwidth}{@{} l L @{}}
\caption[Phase 1 Data Preparation Configuration]{Phase 1 Data Preparation Configuration --- Input, Process, and Output}\label{tab:phase1_config} \\
\toprule
\thead{Dimension} & \thead{Details} \\
\midrule
\endfirsthead
\multicolumn{2}{@{}l}{\textit{(\,Tab.~\ref{tab:phase1_config} continued from previous page\,)}} \\
\toprule
\thead{Dimension} & \thead{Details} \\
\midrule
\endhead
\bottomrule
\endlastfoot
\tablestripes
Objective & Prepare clean, reliable, analysis-ready datasets for all 7 disease domains \\
Focus & Coding, cleaning, screening, bias checks, normality assumptions \\
Input & Raw EEG recordings (EDF), CSV feature matrices, imaging data (DICOM/PNG) \\
Output & Clean validated datasets with documented quality metrics \\
Tools & Python (pandas, scipy, scikit-learn), MNE-Python for EEG \\
Quality Gate & $<$5\% missing values, outliers treated, normality assessed \\
\end{xltabular}}

\subsection{Per-Disease Data Preparation Summary}

Table~\ref{tab:phase1_results} reports the cleaning, screening, and validation metrics for each of the seven disease domains, including sample sizes before and after screening, missing value rates, and imputation methods applied.
\needspace{3\baselineskip}
{\tablestripes\tablefontsize
\begin{xltabular}{\textwidth}{@{} >{\raggedright\arraybackslash}p{2.8cm} R R R R R R R @{}}
\caption[Phase 1 Data Preparation Results by Disease]{Phase 1 Data Preparation Results --- Cleaning, Screening, and Validation Metrics per Disease Domain}\label{tab:phase1_results} \\
\toprule
\thead{Metric} & \thead{SZ} & \thead{EP} & \thead{ST} & \thead{ASD} & \thead{PD} & \thead{DEP} & \thead{CA} \\
\midrule
\endfirsthead
\multicolumn{8}{@{}l}{\textit{(\,Tab.~\ref{tab:phase1_results} continued from previous page\,)}} \\
\toprule
\thead{Metric} & \thead{SZ} & \thead{EP} & \thead{ST} & \thead{ASD} & \thead{PD} & \thead{DEP} & \thead{CA} \\
\midrule
\endhead
\bottomrule
\endlastfoot
\tablestripes
Raw samples & 84 & 102 & 120 & 300 & 50 & 112 & 20{,}000 \\
Post-screen samples & 81 & 98 & 116 & 289 & 48 & 108 & 19{,}847 \\
Exclusion rate (\%) & 3.6 & 3.9 & 3.3 & 3.7 & 4.0 & 3.6 & 0.8 \\
Missing values (\%) & 1.2 & 0.8 & 2.1 & 1.5 & 0.6 & 1.8 & 0.3 \\
Imputation method & KNN & KNN & KNN & KNN & KNN & KNN & Median \\
Duplicates removed & 2 & 3 & 1 & 8 & 1 & 3 & 142 \\
Outliers (IQR, \%) & 4.8 & 3.2 & 5.1 & 3.9 & 2.8 & 4.5 & 1.2 \\
Outlier treatment & Winsor. & Winsor. & Winsor. & Winsor. & Winsor. & Winsor. & Winsor. \\
Normality & Non-N & Non-N & Non-N & Non-N & Non-N & Non-N & Non-N \\
Final feature count & 42 & 38 & 45 & 36 & 40 & 41 & 28 \\
\end{xltabular}}
\vspace{0.5em}\noindent\textit{Note.} SZ = Schizophrenia, EP = Epilepsy, ST = Stress, ASD = Autism Spectrum Disorder, PD = Parkinson's Disease, DEP = Depression, CA = Cancer. Winsor.\ = Winsorised (1st--99th percentile); Non-N = Non-normal (3+/4 normality tests rejecting $H_0$ at $\alpha=0.05$). All EEG domains used KNN imputation ($k=5$, distance-weighted). Cancer used median imputation due to MNAR mechanism in radiomics features.

\subsection{Data Preparation Pipeline}
\needspace{20\baselineskip}
\begin{figure}[!htbp]
\centering
\resizebox{0.97\textwidth}{!}{%
\begin{tikzpicture}[
    pstep/.style={rectangle, rounded corners=3pt, draw=ggunavy, fill=ggunavy!6, minimum width=2.6cm, minimum height=0.9cm, font=\small\bfseries, text=ggunavy, align=center},
    check/.style={diamond, draw=ggunavy, fill=yellow!10, minimum width=1.2cm, minimum height=1.0cm, font=\small, text=ggunavy, align=center, aspect=2},
    arr/.style={-{Stealth[length=4mm, width=3mm]}, very thick, ggunavy}
]
\node[pstep, fill=blue!10] (s1) at (0,0) {Step 1\\Import \& Freeze};
\node[pstep, fill=cyan!10] (s2) at (3.2,0) {Step 2\\Codebook};
\node[pstep, fill=green!10] (s3) at (6.4,0) {Step 3\\Type Validation};
\node[pstep, fill=yellow!12] (s4) at (9.6,0) {Step 4\\Screening};
\node[pstep, fill=orange!12] (s5) at (12.8,0) {Step 5\\Missing Values};

\node[pstep, fill=pink!15] (s6) at (0,-2) {Step 6\\Outlier Detection};
\node[pstep, fill=violet!10] (s7) at (3.2,-2) {Step 7\\Normality Test};
\node[pstep, fill=teal!12] (s8) at (6.4,-2) {Step 8\\Transformation};
\node[pstep, fill=lime!12] (s9) at (9.6,-2) {Step 9\\Bias Check};
\node[pstep, fill=red!10] (s10) at (12.8,-2) {Step 10\\Final Dataset};

\draw[arr] (s1) -- (s2);
\draw[arr] (s2) -- (s3);
\draw[arr] (s3) -- (s4);
\draw[arr] (s4) -- (s5);
\draw[arr] (s5.south) -- ++(0,-0.3) -| (s6.north);
\draw[arr] (s6) -- (s7);
\draw[arr] (s7) -- (s8);
\draw[arr] (s8) -- (s9);
\draw[arr] (s9) -- (s10);
\end{tikzpicture}%
}%
\caption[Phase 1 data preparation pipeline]{Phase 1 data preparation pipeline: ten-step process from raw data import through codebook creation, type validation, eligibility screening, missing value treatment, outlier detection, normality testing, data transformation, bias checking, and final validated dataset production.}
\label{fig:phase1_pipeline}
\end{figure}

\subsection{Missing Value Analysis by Disease}

Table~\ref{tab:phase1_missing} summarises the missing value prevalence, missingness mechanism classification, and imputation strategy applied to each disease domain.
\needspace{8\baselineskip}
{\tablestripes\tablefontsize
\begin{xltabular}{\textwidth}{@{} l R l l R @{}}
\caption[Missing Value Analysis per Disease Domain]{Missing Value Analysis --- Prevalence, Mechanism, and Imputation per Disease Domain}\label{tab:phase1_missing} \\
\toprule
\thead{Disease} & \thead{Missing (\%)} & \thead{Mechanism} & \thead{Imputation} & \thead{Post-Imputation (\%)} \\
\midrule
\endfirsthead
\multicolumn{5}{@{}l}{\textit{(\,Tab.~\ref{tab:phase1_missing} continued from previous page\,)}} \\
\toprule
\thead{Disease} & \thead{Missing (\%)} & \thead{Mechanism} & \thead{Imputation} & \thead{Post-Imputation (\%)} \\
\midrule
\endhead
\bottomrule
\endlastfoot
\tablestripes
Schizophrenia & 1.2 & MCAR & KNN ($k$=5) & 0.0 \\
Epilepsy & 0.8 & MCAR & KNN ($k$=5) & 0.0 \\
Stress & 2.1 & MAR & KNN ($k$=5) & 0.0 \\
ASD & 1.5 & MCAR & KNN ($k$=5) & 0.0 \\
Parkinson & 0.6 & MCAR & KNN ($k$=5) & 0.0 \\
Depression & 1.8 & MAR & KNN ($k$=5) & 0.0 \\
Cancer & 0.3 & MNAR & Median & 0.0 \\
\end{xltabular}}
\vspace{0.5em}\noindent\textit{Note.} MCAR = Missing Completely At Random, MAR = Missing At Random, MNAR = Missing Not At Random. Mechanism classified using correlation of missingness indicators with observed values (MCAR: $r < 0.1$, MAR: $0.1 \leq r < 0.3$, MNAR: $r \geq 0.3$).

\subsection{Outlier Detection Results by Disease}

Table~\ref{tab:phase1_outliers} compares outlier detection results from three methods (IQR, Z-score, Mahalanobis distance) and reports the consensus rate and treatment applied per disease domain.
\needspace{8\baselineskip}
{\tablestripes\tablefontsize
\begin{xltabular}{\textwidth}{@{} l R R R R l @{}}
\caption[Outlier Detection Results per Disease Domain]{Outlier Detection Results --- Three-Method Consensus per Disease Domain}\label{tab:phase1_outliers} \\
\toprule
\thead{Disease} & \thead{IQR (\%)} & \thead{Z-score (\%)} & \thead{Mahalanobis (\%)} & \thead{Consensus (\%)} & \thead{Treatment} \\
\midrule
\endfirsthead
\multicolumn{6}{@{}l}{\textit{(\,Tab.~\ref{tab:phase1_outliers} continued from previous page\,)}} \\
\toprule
\thead{Disease} & \thead{IQR (\%)} & \thead{Z-score (\%)} & \thead{Mahalanobis (\%)} & \thead{Consensus (\%)} & \thead{Treatment} \\
\midrule
\endhead
\bottomrule
\endlastfoot
\tablestripes
Schizophrenia & 4.8 & 3.1 & 2.9 & 2.5 & Winsorise 1--99\textsuperscript{th} \\
Epilepsy & 3.2 & 2.4 & 1.8 & 1.6 & Winsorise 1--99\textsuperscript{th} \\
Stress & 5.1 & 3.8 & 3.2 & 2.9 & Winsorise 1--99\textsuperscript{th} \\
ASD & 3.9 & 2.7 & 2.1 & 1.8 & Winsorise 1--99\textsuperscript{th} \\
Parkinson & 2.8 & 1.9 & 1.5 & 1.2 & Winsorise 1--99\textsuperscript{th} \\
Depression & 4.5 & 3.3 & 2.7 & 2.3 & Winsorise 1--99\textsuperscript{th} \\
Cancer & 1.2 & 0.8 & 0.6 & 0.4 & Winsorise 1--99\textsuperscript{th} \\
\end{xltabular}}
\vspace{0.5em}\noindent\textit{Note.} Consensus = flagged by $\geq$2 of 3 methods. Mahalanobis threshold: $\chi^2_{0.999}(p)$ where $p$ = number of features. All domains treated by winsorisation (capping at 1st and 99th percentile) to preserve sample size.

\subsection{Normality Testing Results by Disease}

Table~\ref{tab:phase1_normality} presents the four-method normality assessment (Shapiro--Wilk, D'Agostino, Kolmogorov--Smirnov, Anderson--Darling) and the consensus verdict for each disease domain.
\needspace{8\baselineskip}
{\tablestripes\tablefontsize
\begin{xltabular}{\textwidth}{@{} >{\raggedright\arraybackslash}p{2.2cm} R R R L L @{}}
\caption[Normality Test Results per Disease Domain]{Normality Test Results --- Four-Method Assessment per Disease Domain}\label{tab:phase1_normality} \\
\toprule
\thead{Disease} & \thead{Shapiro-Wilk} & \thead{D'Agostino} & \thead{KS Test} & \thead{Anderson-Darling} & \thead{Consensus} \\
\midrule
\endfirsthead
\multicolumn{6}{@{}l}{\textit{(\,Tab.~\ref{tab:phase1_normality} continued from previous page\,)}} \\
\toprule
\thead{Disease} & \thead{Shapiro-Wilk} & \thead{D'Agostino} & \thead{KS Test} & \thead{Anderson-Darling} & \thead{Consensus} \\
\midrule
\endhead
\bottomrule
\endlastfoot
\tablestripes
Schizophrenia & $p<0.001$ & $p<0.001$ & $p=0.003$ & Reject & Non-normal (4/4) \\
Epilepsy & $p<0.001$ & $p<0.001$ & $p=0.008$ & Reject & Non-normal (4/4) \\
Stress & $p<0.001$ & $p<0.001$ & $p=0.012$ & Reject & Non-normal (4/4) \\
ASD & $p<0.001$ & $p<0.001$ & $p=0.005$ & Reject & Non-normal (4/4) \\
Parkinson & $p=0.002$ & $p=0.004$ & $p=0.031$ & Reject & Non-normal (4/4) \\
Depression & $p<0.001$ & $p<0.001$ & $p=0.009$ & Reject & Non-normal (4/4) \\
Cancer & $p<0.001$ & $p<0.001$ & $p<0.001$ & Reject & Non-normal (4/4) \\
\end{xltabular}}
\vspace{0.5em}\noindent\textit{Note.} All seven domains exhibit non-normal distributions, justifying the use of non-parametric tests (Mann-Whitney U, Wilcoxon signed-rank) and robust estimation methods (bootstrap resampling with 1,000 resamples) throughout the analysis.

\subsection{Phase 1 Decision Points}

Table~\ref{tab:phase1_decisions} enumerates the six key decision points encountered during data preparation, together with the thresholds applied and actions taken at each gate.
\needspace{8\baselineskip}
{\tablestripes\tablefontsize
\begin{xltabular}{\textwidth}{@{} l l l L @{}}
\caption[Phase 1 Decision Points and Actions]{Phase 1 Decision Points --- Conditions, Thresholds, and Actions Taken}\label{tab:phase1_decisions} \\
\toprule
\thead{Decision} & \thead{Condition} & \thead{Threshold} & \thead{Action Taken} \\
\midrule
\endfirsthead
\multicolumn{4}{@{}l}{\textit{(\,Tab.~\ref{tab:phase1_decisions} continued from previous page\,)}} \\
\toprule
\thead{Decision} & \thead{Condition} & \thead{Threshold} & \thead{Action Taken} \\
\midrule
\endhead
\bottomrule
\endlastfoot
\tablestripes
Missing data severity & $>$10\% missing & 10\% & No domain exceeded 2.1\%; KNN imputation applied \\
Imputation method & MCAR vs MAR vs MNAR & Correlation $r$ & KNN for MCAR/MAR; median for MNAR  \\
Outlier treatment & Consensus $\geq$2/3 methods & 2+ methods & Winsorisation (1st--99th) to preserve sample size \\
Normality assumption & 3+/4 tests reject $H_0$ & $\alpha=0.05$ & Non-parametric tests selected for all domains \\
Duplicate removal & Exact row matches & 100\% match & 160 total duplicates removed for ASD \\
Sample adequacy & Minimum per domain & $n \geq 30$ & All domains exceed minimum (smallest: PD $n=48$) \\
\end{xltabular}}

% [SPACE-FIX] clearpage removed to eliminate empty page

%% ======================================================================
%% G.2  PHASE 2 — DESCRIPTIVE ANALYSIS (Per Disease)
%% ======================================================================
\section{Phase 2 --- Descriptive Analysis}
\label{sec:phase2_descriptive}

Phase 2 summarises sample characteristics and variable behaviour for each disease domain before proceeding to reliability testing and factor analysis.

\subsection{Phase 2 Master Configuration}

Table~\ref{tab:phase2_config} documents the master configuration for Phase~2, specifying the objectives, focus areas, inputs, outputs, and quality gates that govern descriptive analysis.
\needspace{8\baselineskip}
{\tablestripes\tablefontsize
\begin{xltabular}{\textwidth}{@{} l L @{}}
\caption[Phase 2 Descriptive Analysis Configuration]{Phase 2 Descriptive Analysis Configuration --- Objectives and Methods}\label{tab:phase2_config} \\
\toprule
\thead{Dimension} & \thead{Details} \\
\midrule
\endfirsthead
\multicolumn{2}{@{}l}{\textit{(\,Tab.~\ref{tab:phase2_config} continued from previous page\,)}} \\
\toprule
\thead{Dimension} & \thead{Details} \\
\midrule
\endhead
\bottomrule
\endlastfoot
\tablestripes
Objective & Summarise sample profiles and variable distributions across all 7 disease domains \\
Focus & Frequencies, central tendency, dispersion, distribution shape \\
Input & Clean validated datasets from Phase 1 \\
Output & Sample profiles, descriptive statistics tables, initial construct insights \\
Tools & Python (pandas, numpy, scipy.stats) \\
Quality Gate & All variables summarised; no unexplained anomalies \\
\end{xltabular}}

\subsection{Sample Profile by Disease}

Table~\ref{tab:phase2_profile} maps the sample size, age range, sex distribution, data source, file format, and sampling rate for each of the seven disease domains.
\needspace{8\baselineskip}
{\tablestripes\tablefontsize
\begin{xltabular}{\textwidth}{@{} l R l R l l l @{}}
\caption[Sample Profile Summary by Disease Domain]{Sample Profile Summary --- Subject Demographics and Recording Parameters per Disease Domain}\label{tab:phase2_profile} \\
\toprule
\thead{Disease} & \thead{$n$} & \thead{Age Range} & \thead{Female (\%)} & \thead{Source} & \thead{Format} & \thead{Hz} \\
\midrule
\endfirsthead
\multicolumn{7}{@{}l}{\textit{(\,Tab.~\ref{tab:phase2_profile} continued from previous page\,)}} \\
\toprule
\thead{Disease} & \thead{$n$} & \thead{Age Range} & \thead{Female (\%)} & \thead{Source} & \thead{Format} & \thead{Hz} \\
\midrule
\endhead
\bottomrule
\endlastfoot
\tablestripes
Schizophrenia & 81 & 18--60 & 41.2 & RepOD & EDF & 128 \\
Epilepsy & 98 & 5--55 & 48.0 & CHB-MIT & EDF & 256 \\
Stress & 116 & 20--45 & 52.6 & DEAP/SAM40 & CSV & 512/128 \\
ASD & 289 & 3--18 & 22.5 & KAU/Kaggle & CSV & 128 \\
Parkinson & 48 & 50--80 & 35.4 & PhysioNet & EDF & 256 \\
Depression & 108 & 18--65 & 55.6 & OpenNeuro & EDF & 256 \\
Cancer & 19{,}847 & 25--85 & 51.3 & TCIA & DICOM & --- \\
\end{xltabular}}

\subsection{Descriptive Statistics by Disease}

Table~\ref{tab:phase2_descriptive} presents the central tendency, dispersion, and distribution shape statistics aggregated across all extracted features for each disease domain.
\needspace{3\baselineskip}
{\tablestripes\tablefontsize
\begin{xltabular}{\textwidth}{@{} l R R R R R R R @{}}
\caption[Descriptive Statistics Summary Across All Disease Domains]{Descriptive Statistics Summary --- Central Tendency and Dispersion of Key Features per Disease Domain}\label{tab:phase2_descriptive} \\
\toprule
\thead{Statistic} & \thead{SZ} & \thead{EP} & \thead{ST} & \thead{ASD} & \thead{PD} & \thead{DEP} & \thead{CA} \\
\midrule
\endfirsthead
\multicolumn{8}{@{}l}{\textit{(\,Tab.~\ref{tab:phase2_descriptive} continued from previous page\,)}} \\
\toprule
\thead{Statistic} & \thead{SZ} & \thead{EP} & \thead{ST} & \thead{ASD} & \thead{PD} & \thead{DEP} & \thead{CA} \\
\midrule
\endhead
\bottomrule
\endlastfoot
\tablestripes
Mean ($\mu$) & 3.42 & 4.18 & 3.95 & 4.12 & 3.78 & 3.65 & 0.52 \\
Median & 3.28 & 4.05 & 3.82 & 4.01 & 3.65 & 3.51 & 0.48 \\
Std. deviation ($\sigma$) & 1.45 & 1.82 & 1.38 & 1.21 & 1.56 & 1.41 & 0.18 \\
Variance ($\sigma^2$) & 2.10 & 3.31 & 1.90 & 1.46 & 2.43 & 1.99 & 0.03 \\
Range & 8.74 & 12.31 & 7.62 & 6.85 & 9.18 & 8.92 & 0.95 \\
IQR & 1.82 & 2.34 & 1.71 & 1.48 & 1.95 & 1.78 & 0.22 \\
Skewness & 0.85 & 1.12 & 0.72 & 0.58 & 0.94 & 0.81 & 0.34 \\
Kurtosis & 2.41 & 3.82 & 1.95 & 1.42 & 2.78 & 2.31 & $-0.12$ \\
CoV (\%) & 42.4 & 43.5 & 34.9 & 29.4 & 41.3 & 38.6 & 34.6 \\
\end{xltabular}}
\vspace{0.5em}\noindent\textit{Note.} Values represent aggregated statistics across all extracted features per domain. CoV = Coefficient of Variation. Cancer features are normalised to $[0,1]$ range (radiomics), explaining the smaller absolute values.

\subsection{Distribution Shape Visualisation}
\needspace{20\baselineskip}
\begin{figure}[!htbp]
\centering
\resizebox{0.97\textwidth}{!}{%
\begin{tikzpicture}[
    bar/.style={fill=ggunavy!60, draw=ggunavy},
    font=\small
]
%% Skewness comparison bars
\draw[-{Stealth[length=4mm,width=3mm]}, thick, ggunavy!50] (0,0) -- (0,4.5) node[above, font=\small] {Skewness};
\draw[-{Stealth[length=4mm,width=3mm]}, thick, ggunavy!50] (0,0) -- (11,0);

%% Disease bars
\fill[bar, fill=blue!40] (0.5,0) rectangle (1.2,0.85*3);
\fill[bar, fill=cyan!40] (1.8,0) rectangle (2.5,1.12*3);
\fill[bar, fill=green!40] (3.1,0) rectangle (3.8,0.72*3);
\fill[bar, fill=yellow!50] (4.4,0) rectangle (5.1,0.58*3);
\fill[bar, fill=orange!40] (5.7,0) rectangle (6.4,0.94*3);
\fill[bar, fill=pink!50] (7.0,0) rectangle (7.7,0.81*3);
\fill[bar, fill=violet!30] (8.3,0) rectangle (9.0,0.34*3);

%% Labels
\node[rotate=30, anchor=east] at (0.85,-0.3) {SZ};
\node[rotate=30, anchor=east] at (2.15,-0.3) {EP};
\node[rotate=30, anchor=east] at (3.45,-0.3) {ST};
\node[rotate=30, anchor=east] at (4.75,-0.3) {ASD};
\node[rotate=30, anchor=east] at (6.05,-0.3) {PD};
\node[rotate=30, anchor=east] at (7.35,-0.3) {DEP};
\node[rotate=30, anchor=east] at (8.65,-0.3) {CA};

%% Values
\node[above, font=\small] at (0.85,0.85*3) {0.85};
\node[above, font=\small] at (2.15,1.12*3) {1.12};
\node[above, font=\small] at (3.45,0.72*3) {0.72};
\node[above, font=\small] at (4.75,0.58*3) {0.58};
\node[above, font=\small] at (6.05,0.94*3) {0.94};
\node[above, font=\small] at (7.35,0.81*3) {0.81};
\node[above, font=\small] at (8.65,0.34*3) {0.34};

%% Y-axis ticks
\foreach \y/\l in {0/0, 1.5/0.5, 3/1.0} {
    \draw[ggunavy!30] (-0.1,\y) -- (10,\y);
    \node[left, font=\small] at (-0.1,\y) {\l};
}

%% Threshold line
\draw[red!50, dashed, thick] (0,3) -- (10,3);
\node[right, font=\small, text=red!60] at (10,3) {|skew|=1};
\end{tikzpicture}%
}%
\caption[Skewness comparison across disease domains]{Skewness comparison across all seven disease domains. All domains show positive skewness (right-tailed distribution). Epilepsy exceeds the $|$skewness$|>1$ threshold, indicating moderate right skew from high-amplitude ictal events. Cancer shows near-symmetric distribution due to radiomics normalisation.}
\label{fig:phase2_skewness}
\end{figure}

\subsection{Construct-Level Summary}

% MOVED TO CHAPTER 4 MAIN TEXT (Mar 2026)
\iffalse
\needspace{8\baselineskip}
\noindent Table~\ref{tab:app_phase2_construct} summarises construct-level descriptive summary --- feature category means and variability per disease for appendix review.

{\tablestripes\tablefontsize
\begin{xltabular}{\textwidth}{@{} l l R R l @{}}
\caption[Construct-Level Descriptive Summary]{Construct-Level Descriptive Summary --- Feature Category Means and Variability per Disease Domain}\label{tab:app_phase2_construct} \\
\toprule
\thead{Disease} & \thead{Feature Category} & \thead{Mean} & \thead{SD} & \thead{Interpretation} \\
\midrule
\endfirsthead
\multicolumn{5}{@{}l}{\textit{(\,Tab.~\ref{tab:app_phase2_construct} continued from previous page\,)}} \\
\toprule
\thead{Disease} & \thead{Feature Category} & \thead{Mean} & \thead{SD} & \thead{Interpretation} \\
\midrule
\endhead
\bottomrule
\endlastfoot
\tablestripes
SZ & Alpha power & 2.18 & 0.95 & Reduced (deficit signature) \\
SZ & Theta/beta ratio & 3.42 & 1.12 & Elevated (attentional deficit) \\
EP & Delta power & 6.82 & 2.45 & Highly elevated (ictal signature) \\
EP & Line length & 5.91 & 1.78 & Elevated (seizure marker) \\
ST & EDA (SCL) & 4.25 & 1.38 & Elevated (arousal marker) \\
ST & HRV (RMSSD) & 32.4 & 12.8 & Reduced (stress response) \\
ASD & F-P coherence & 0.42 & 0.18 & Reduced (connectivity deficit) \\
ASD & Gamma power & 1.85 & 0.72 & Reduced (binding deficit) \\
PD & Beta power & 1.92 & 0.88 & Suppressed (motor deficit) \\
PD & Tremor frequency & 4.8 & 1.2 & Characteristic 4--6 Hz band \\
DEP & Alpha asymmetry & $-0.32$ & 0.21 & Left-lateralised deficit \\
DEP & Theta power & 3.78 & 1.15 & Elevated (rumination marker) \\
CA & GLCM contrast & 0.68 & 0.15 & Elevated (tumour texture) \\
CA & Shape irregularity & 0.74 & 0.12 & Elevated (malignancy marker) \\
\end{xltabular}}
\fi

% [SPACE-FIX] clearpage removed to eliminate empty page

%% ======================================================================
%% G.3  PHASE 3 — RELIABILITY & VALIDITY (Per Disease)
%% ======================================================================
\section{Phase 3 --- Reliability and Validity}
\label{sec:phase3_reliability}

Phase 3 validates measurement quality of constructs within each disease domain, ensuring internal consistency and construct validity before factor analysis.

\subsection{Internal Consistency (Cronbach's Alpha)}

Table~\ref{tab:phase3_alpha} reports Cronbach's alpha values for each feature category within every disease domain, confirming that all constructs meet the 0.70 reliability threshold.
\needspace{8\baselineskip}
{\tablestripes\tablefontsize
\begin{xltabular}{\textwidth}{@{} l R R R R R R R @{}}
\caption[Cronbach's Alpha by Feature Category and Disease]{Internal Consistency (Cronbach's $\alpha$) --- Feature Category Reliability per Disease Domain}\label{tab:phase3_alpha} \\
\toprule
\thead{Feature Category} & \thead{SZ} & \thead{EP} & \thead{ST} & \thead{ASD} & \thead{PD} & \thead{DEP} & \thead{CA} \\
\midrule
\endfirsthead
\multicolumn{8}{@{}l}{\textit{(\,Tab.~\ref{tab:phase3_alpha} continued from previous page\,)}} \\
\toprule
\thead{Feature Category} & \thead{SZ} & \thead{EP} & \thead{ST} & \thead{ASD} & \thead{PD} & \thead{DEP} & \thead{CA} \\
\midrule
\endhead
\bottomrule
\endlastfoot
\tablestripes
Time-domain & 0.89 & 0.91 & 0.87 & 0.92 & 0.88 & 0.90 & --- \\
Frequency-domain & 0.92 & 0.93 & 0.90 & 0.91 & 0.89 & 0.91 & --- \\
Connectivity & 0.85 & 0.82 & 0.81 & 0.88 & 0.84 & 0.86 & --- \\
Spatial & 0.87 & 0.84 & 0.83 & 0.86 & 0.85 & 0.87 & --- \\
Radiomics (1st order) & --- & --- & --- & --- & --- & --- & 0.94 \\
Radiomics (texture) & --- & --- & --- & --- & --- & --- & 0.91 \\
Radiomics (shape) & --- & --- & --- & --- & --- & --- & 0.88 \\
Wearable (ECG/EDA) & --- & --- & 0.86 & --- & --- & --- & --- \\
\end{xltabular}}
\vspace{0.5em}\noindent\textit{Note.} All $\alpha$ values exceed the 0.70 threshold, indicating good to excellent internal consistency. ``---'' indicates feature category not applicable to that domain.

\subsection{Composite Reliability and Convergent Validity}

Table~\ref{tab:phase3_cr_ave} presents composite reliability and average variance extracted values per feature category and disease domain, establishing convergent validity across all constructs.
\needspace{8\baselineskip}
{\tablestripes\tablefontsize
\begin{xltabular}{\textwidth}{@{} l R R R R R R R @{}}
\caption[Composite Reliability and AVE by Disease Domain]{Composite Reliability (CR) and Average Variance Extracted (AVE) per Disease Domain}\label{tab:phase3_cr_ave} \\
\toprule
\thead{Metric} & \thead{SZ} & \thead{EP} & \thead{ST} & \thead{ASD} & \thead{PD} & \thead{DEP} & \thead{CA} \\
\midrule
\endfirsthead
\multicolumn{8}{@{}l}{\textit{(\,Tab.~\ref{tab:phase3_cr_ave} continued from previous page\,)}} \\
\toprule
\thead{Metric} & \thead{SZ} & \thead{EP} & \thead{ST} & \thead{ASD} & \thead{PD} & \thead{DEP} & \thead{CA} \\
\midrule
\endhead
\bottomrule
\endlastfoot
\tablestripes
CR (Time) & 0.91 & 0.93 & 0.89 & 0.94 & 0.90 & 0.92 & --- \\
CR (Frequency) & 0.94 & 0.95 & 0.92 & 0.93 & 0.91 & 0.93 & --- \\
CR (Connectivity) & 0.88 & 0.85 & 0.84 & 0.90 & 0.87 & 0.89 & --- \\
CR (Radiomics) & --- & --- & --- & --- & --- & --- & 0.95 \\
AVE (Time) & 0.62 & 0.68 & 0.58 & 0.72 & 0.61 & 0.65 & --- \\
AVE (Frequency) & 0.71 & 0.74 & 0.65 & 0.70 & 0.64 & 0.69 & --- \\
AVE (Connectivity) & 0.55 & 0.52 & 0.51 & 0.60 & 0.54 & 0.57 & --- \\
AVE (Radiomics) & --- & --- & --- & --- & --- & --- & 0.73 \\
\end{xltabular}}
\vspace{0.5em}\noindent\textit{Note.} CR threshold: $>$0.70 (all satisfied). AVE threshold: $>$0.50 (all satisfied). These values confirm convergent validity across all feature categories and disease domains.

\subsection{Discriminant Validity (HTMT)}

Table~\ref{tab:phase3_htmt} documents the Heterotrait--Monotrait ratio values for all feature category pairs across six EEG disease domains, verifying that all pairs fall below the 0.85 strict threshold.
\needspace{8\baselineskip}
{\tablestripes\tablefontsize
\begin{xltabular}{\textwidth}{@{} l R R R R R R @{}}
\caption[HTMT Discriminant Validity by Disease Domain]{Discriminant Validity (HTMT) --- Feature Category Pair Values per Disease Domain}\label{tab:phase3_htmt} \\
\toprule
\thead{Feature Pair} & \thead{SZ} & \thead{EP} & \thead{ST} & \thead{ASD} & \thead{PD} & \thead{DEP} \\
\midrule
\endfirsthead
\multicolumn{7}{@{}l}{\textit{(\,Tab.~\ref{tab:phase3_htmt} continued from previous page\,)}} \\
\toprule
\thead{Feature Pair} & \thead{SZ} & \thead{EP} & \thead{ST} & \thead{ASD} & \thead{PD} & \thead{DEP} \\
\midrule
\endhead
\bottomrule
\endlastfoot
\tablestripes
Time--Frequency & 0.72 & 0.68 & 0.71 & 0.65 & 0.74 & 0.69 \\
Time--Connectivity & 0.58 & 0.54 & 0.52 & 0.61 & 0.56 & 0.55 \\
Frequency--Connectivity & 0.64 & 0.61 & 0.59 & 0.67 & 0.62 & 0.63 \\
Time--Spatial & 0.51 & 0.48 & 0.47 & 0.53 & 0.50 & 0.49 \\
\end{xltabular}}
\vspace{0.5em}\noindent\textit{Note.} All HTMT values $<$0.85 (strict threshold), confirming discriminant validity. Feature categories are sufficiently distinct from each other within each disease domain.

\subsection{Phase 3 Quality Assessment}
\needspace{20\baselineskip}
\begin{figure}[!htbp]
\centering
\resizebox{0.97\textwidth}{!}{%
\begin{tikzpicture}[font=\small]
%% Axes
\draw[-{Stealth[length=4mm,width=3mm]}, thick, ggunavy!50] (0,0) -- (0,5) node[above, font=\small] {Score};
\draw[-{Stealth[length=4mm,width=3mm]}, thick, ggunavy!50] (0,0) -- (12,0);

%% Three grouped bars per disease (Alpha, CR, AVE)
\foreach \x/\lab/\a/\c/\v in {
    1/SZ/0.88/0.91/0.63,
    3/EP/0.88/0.91/0.65,
    5/ST/0.85/0.88/0.58,
    7/ASD/0.89/0.92/0.67,
    9/PD/0.87/0.89/0.60,
    11/DEP/0.88/0.91/0.64} {
    \fill[blue!40] (\x-0.35,0) rectangle (\x,\a*5);
    \fill[green!40] (\x,0) rectangle (\x+0.35,\c*5);
    \fill[orange!40] (\x+0.35,0) rectangle (\x+0.7,\v*5);
    \node[rotate=30, anchor=east] at (\x+0.17,-0.3) {\lab};
}

%% Threshold lines
\draw[red!50, dashed] (0,0.70*5) -- (12,0.70*5);
\node[right, font=\small, text=red!50] at (12,0.70*5) {CR/\,$\alpha$=0.70};
\draw[red!50, dotted] (0,0.50*5) -- (12,0.50*5);
\node[right, font=\small, text=red!50] at (12,0.50*5) {AVE=0.50};

%% Y ticks
\foreach \y/\l in {0/0, 1.25/0.25, 2.5/0.50, 3.75/0.75, 5/1.00} {
    \node[left, font=\small] at (-0.1,\y) {\l};
}

%% Legend
\fill[blue!40] (0.5,4.5) rectangle (1,4.7);
\node[right, font=\small] at (1,4.6) {$\alpha$};
\fill[green!40] (2.5,4.5) rectangle (3,4.7);
\node[right, font=\small] at (3,4.6) {CR};
\fill[orange!40] (4.5,4.5) rectangle (5,4.7);
\node[right, font=\small] at (5,4.6) {AVE};

\end{tikzpicture}%
}%
\caption[Reliability and validity comparison across diseases]{Reliability and validity comparison across six EEG disease domains. All Cronbach's $\alpha$ and CR values exceed 0.70; all AVE values exceed 0.50. Cancer domain (not shown) achieves $\alpha=0.91$, CR=0.95, AVE=0.73 for radiomics features.}
\label{fig:phase3_reliability}
\end{figure}

% [SPACE-FIX] clearpage removed to eliminate empty page

%% ======================================================================
%% G.4  PHASE 4 — FACTOR ANALYSIS (Per Disease)
%% ======================================================================
\section{Phase 4 --- Factor Analysis (EFA + CFA)}
\label{sec:phase4_factor}

Phase 4 identifies latent factor structures through exploratory factor analysis (EFA) and confirms them through confirmatory factor analysis (CFA) for each disease domain.

\subsection{KMO and Bartlett's Test}

Table~\ref{tab:phase4_kmo} reports the Kaiser--Meyer--Olkin measure of sampling adequacy and Bartlett's test of sphericity for each disease domain, confirming suitability for factor analysis.
\needspace{8\baselineskip}
{\tablestripes\tablefontsize
\begin{xltabular}{\textwidth}{@{} l R R l @{}}
\caption[KMO and Bartlett's Test Results per Disease]{KMO and Bartlett's Test --- Factor Analysis Suitability per Disease Domain}\label{tab:phase4_kmo} \\
\toprule
\thead{Disease} & \thead{KMO} & \thead{Bartlett's $\chi^2$} & \thead{Suitability} \\
\midrule
\endfirsthead
\multicolumn{4}{@{}l}{\textit{(\,Tab.~\ref{tab:phase4_kmo} continued from previous page\,)}} \\
\toprule
\thead{Disease} & \thead{KMO} & \thead{Bartlett's $\chi^2$} & \thead{Suitability} \\
\midrule
\endhead
\bottomrule
\endlastfoot
\tablestripes
Schizophrenia & 0.84 & 1,245.3 ($p<0.001$) & Good \\
Epilepsy & 0.87 & 1,892.7 ($p<0.001$) & Meritorious \\
Stress & 0.82 & 1,078.5 ($p<0.001$) & Good \\
ASD & 0.89 & 3,412.8 ($p<0.001$) & Meritorious \\
Parkinson & 0.78 & 628.4 ($p<0.001$) & Good \\
Depression & 0.85 & 1,521.6 ($p<0.001$) & Good \\
Cancer & 0.91 & 42,185.2 ($p<0.001$) & Marvellous \\
\end{xltabular}}
\vspace{0.5em}\noindent\textit{Note.} KMO thresholds: $>$0.90 marvellous, $>$0.80 meritorious/good, $>$0.70 middling, $<$0.50 unacceptable. All domains exceed the 0.70 minimum, confirming suitability for factor analysis.

\subsection{EFA Factor Extraction}

Table~\ref{tab:phase4_efa} summarises the exploratory factor analysis results, including the number of factors retained, cumulative variance explained, eigenvalue cutoff, and rotation method for each disease domain.
\needspace{8\baselineskip}
{\tablestripes\tablefontsize
\begin{xltabular}{\textwidth}{@{} l R R R l @{}}
\caption[EFA Factor Extraction Results per Disease]{EFA Factor Extraction --- Eigenvalues, Factors Retained, and Variance Explained per Disease Domain}\label{tab:phase4_efa} \\
\toprule
\thead{Disease} & \thead{Factors Retained} & \thead{Variance (\%)} & \thead{Eigenvalue Cutoff} & \thead{Rotation} \\
\midrule
\endfirsthead
\multicolumn{5}{@{}l}{\textit{(\,Tab.~\ref{tab:phase4_efa} continued from previous page\,)}} \\
\toprule
\thead{Disease} & \thead{Factors Retained} & \thead{Variance (\%)} & \thead{Eigenvalue Cutoff} & \thead{Rotation} \\
\midrule
\endhead
\bottomrule
\endlastfoot
\tablestripes
Schizophrenia & 4 & 72.4 & 1.0 & Promax \\
Epilepsy & 4 & 75.8 & 1.0 & Promax \\
Stress & 5 & 71.2 & 1.0 & Promax \\
ASD & 3 & 78.5 & 1.0 & Promax \\
Parkinson & 4 & 69.8 & 1.0 & Promax \\
Depression & 4 & 73.1 & 1.0 & Promax \\
Cancer & 3 & 82.3 & 1.0 & Promax \\
\end{xltabular}}
\vspace{0.5em}\noindent\textit{Note.} Factors retained based on eigenvalue $>$1 (Kaiser criterion) confirmed by scree plot. All domains exceed 60\% cumulative variance explained. Promax rotation used (oblique) because constructs are theoretically correlated.

\subsection{CFA Model Fit Indices}

Table~\ref{tab:phase4_cfa} compares confirmatory factor analysis fit indices (CFI, RMSEA, SRMR, $\chi^2$/df) across all seven disease domains, with verdicts ranging from acceptable to excellent fit.
\needspace{8\baselineskip}
{\tablestripes\tablefontsize
\begin{xltabular}{\textwidth}{@{} l R R R R l @{}}
\caption[CFA Model Fit Indices per Disease Domain]{CFA Model Fit Indices --- Measurement Model Validation per Disease Domain}\label{tab:phase4_cfa} \\
\toprule
\thead{Disease} & \thead{CFI} & \thead{RMSEA} & \thead{SRMR} & \thead{$\chi^2$/df} & \thead{Verdict} \\
\midrule
\endfirsthead
\multicolumn{6}{@{}l}{\textit{(\,Tab.~\ref{tab:phase4_cfa} continued from previous page\,)}} \\
\toprule
\thead{Disease} & \thead{CFI} & \thead{RMSEA} & \thead{SRMR} & \thead{$\chi^2$/df} & \thead{Verdict} \\
\midrule
\endhead
\bottomrule
\endlastfoot
\tablestripes
Schizophrenia & 0.94 & 0.058 & 0.042 & 1.82 & Good fit \\
Epilepsy & 0.96 & 0.051 & 0.038 & 1.65 & Excellent fit \\
Stress & 0.93 & 0.062 & 0.048 & 1.94 & Good fit \\
ASD & 0.97 & 0.045 & 0.032 & 1.52 & Excellent fit \\
Parkinson & 0.92 & 0.068 & 0.054 & 2.12 & Acceptable fit \\
Depression & 0.95 & 0.054 & 0.041 & 1.74 & Good fit \\
Cancer & 0.98 & 0.038 & 0.028 & 1.35 & Excellent fit \\
\end{xltabular}}
\vspace{0.5em}\noindent\textit{Note.} Thresholds: CFI $>$0.90, RMSEA $<$0.08, SRMR $<$0.08, $\chi^2$/df $<$3. All seven domains meet acceptable-to-excellent fit criteria, validating the measurement models.

\subsection{Factor Structure Visualisation}
\needspace{20\baselineskip}
\begin{figure}[!htbp]
\centering
\resizebox{0.97\textwidth}{!}{%
\begin{tikzpicture}[
    latent/.style={ellipse, draw=ggunavy, minimum width=2.2cm, minimum height=1cm, font=\small\bfseries, text=ggunavy, align=center},
    obs/.style={rectangle, draw=ggunavy, minimum width=1.5cm, minimum height=1.0cm, font=\small, align=center},
    arr/.style={-{Stealth[length=4mm, width=3mm]}, very thick, ggunavy},
    corr/.style={{Stealth[length=3mm, width=2mm]}-{Stealth[length=3mm, width=2mm]}, thick, ggunavy!40, dashed}
]
%% Latent factors
\node[latent, fill=lightblue] (f1) at (0,0) {Time\\Domain};
\node[latent, fill=lightorange] (f2) at (4.5,0) {Frequency\\Domain};
\node[latent, fill=lightteal] (f3) at (9,0) {Connectivity};
\node[latent, fill=lightpurple] (f4) at (13.5,0) {Spatial};

%% Observed indicators
\foreach \i/\lab in {1/Mean, 2/Var, 3/ZCR, 4/Hjorth} {
    \node[obs, fill=lightblue!40] (o1\i) at ({-1.5+\i*1}, -1.8) {\lab};
    \draw[arr] (f1) -- (o1\i);
}
\foreach \i/\lab in {1/$\delta$, 2/$\theta$, 3/$\alpha$, 4/$\beta$} {
    \node[obs, fill=lightorange!40] (o2\i) at ({3+\i*1}, -1.8) {\lab};
    \draw[arr] (f2) -- (o2\i);
}
\foreach \i/\lab in {1/Coh, 2/PLV, 3/MI} {
    \node[obs, fill=lightteal!40] (o3\i) at ({7.5+\i*1}, -1.8) {\lab};
    \draw[arr] (f3) -- (o3\i);
}
\foreach \i/\lab in {1/Asym, 2/Ratio, 3/Topo} {
    \node[obs, fill=lightpurple!40] (o4\i) at ({12+\i*1}, -1.8) {\lab};
    \draw[arr] (f4) -- (o4\i);
}

%% Correlations between factors
\draw[corr] (f1.north) to[bend left=25] node[above, font=\small] {$r$=0.45} (f2.north);
\draw[corr] (f2.north) to[bend left=25] node[above, font=\small] {$r$=0.38} (f3.north);
\draw[corr] (f3.north) to[bend left=25] node[above, font=\small] {$r$=0.32} (f4.north);
\end{tikzpicture}%
}
\caption[CFA measurement model structure]{CFA measurement model structure showing four latent factors (time-domain, frequency-domain, connectivity, spatial) with their observed indicators. Dashed lines indicate inter-factor correlations. This structure is common across all six EEG disease domains; Cancer uses a separate three-factor model (first-order statistics, texture features, shape descriptors).}
\label{fig:phase4_cfa}
\end{figure}

% [SPACE-FIX] clearpage removed to eliminate empty page

%% ======================================================================
%% G.5  PHASE 5 — REGRESSION + SEM (Per Disease)
%% ======================================================================
\section{Phase 5 --- Dependence Analysis (Regression + SEM)}
\label{sec:phase5_sem}

Phase 5 tests hypothesised relationships between feature constructs and classification outcomes using regression and structural equation modelling (SEM).

\subsection{Regression Results per Disease}

Table~\ref{tab:phase5_regression} presents the standardised regression coefficients for each feature category predictor and the overall model fit ($R^2$, adjusted $R^2$) per disease domain.
\needspace{8\baselineskip}
{\tablestripes\tablefontsize
\begin{xltabular}{\textwidth}{@{} >{\raggedright\arraybackslash}p{2.2cm} R R R R R @{}}
\caption[Multiple Regression Results per Disease Domain]{Multiple Regression Results --- Feature Category Coefficients and Model Fit per Disease Domain}\label{tab:phase5_regression} \\
\toprule
\thead{Disease} & \thead{$\beta$ (Time)} & \thead{$\beta$ (Freq)} & \thead{$\beta$ (Conn)} & \thead{R$^2$} & \thead{Adj.\ R$^2$} \\
\midrule
\endfirsthead
\multicolumn{6}{@{}l}{\textit{(\,Tab.~\ref{tab:phase5_regression} continued from previous page\,)}} \\
\toprule
\thead{Disease} & \thead{$\beta$ (Time)} & \thead{$\beta$ (Freq)} & \thead{$\beta$ (Conn)} & \thead{R$^2$} & \thead{Adj.\ R$^2$} \\
\midrule
\endhead
\bottomrule
\endlastfoot
\tablestripes
Schizophrenia & 0.32* & 0.48** & 0.28* & 0.71 & 0.69 \\
Epilepsy & 0.41** & 0.52** & 0.18 & 0.78 & 0.76 \\
Stress & 0.38** & 0.42** & 0.25* & 0.68 & 0.66 \\
ASD & 0.29* & 0.55** & 0.35** & 0.82 & 0.81 \\
Parkinson & 0.44** & 0.39** & 0.22* & 0.74 & 0.71 \\
Depression & 0.35** & 0.45** & 0.30* & 0.72 & 0.70 \\
Cancer & 0.51** & --- & --- & 0.85 & 0.84 \\
\end{xltabular}}
\vspace{0.5em}\noindent\textit{Note.} * $p<0.05$, ** $p<0.01$. Cancer uses radiomics feature categories (1st-order, texture, shape) instead of EEG categories. Frequency-domain features are the strongest predictors across all EEG domains.

\subsection{SEM Path Coefficients}

Table~\ref{tab:phase5_sem_paths} reports the PLS-SEM structural path coefficients for five hypothesised relationships across all seven disease domains, with bootstrap significance levels.
\needspace{3\baselineskip}
{\tablestripes\tablefontsize
\begin{xltabular}{\textwidth}{@{} >{\raggedright\arraybackslash}p{3.5cm} R R R R R R R @{}}
\caption[SEM Path Coefficients per Disease Domain]{SEM Path Coefficients --- Structural Model Results per Disease Domain}\label{tab:phase5_sem_paths} \\
\toprule
\thead{Path} & \thead{SZ} & \thead{EP} & \thead{ST} & \thead{ASD} & \thead{PD} & \thead{DEP} & \thead{CA} \\
\midrule
\endfirsthead
\multicolumn{8}{@{}l}{\textit{(\,Tab.~\ref{tab:phase5_sem_paths} continued from previous page\,)}} \\
\toprule
\thead{Path} & \thead{SZ} & \thead{EP} & \thead{ST} & \thead{ASD} & \thead{PD} & \thead{DEP} & \thead{CA} \\
\midrule
\endhead
\bottomrule
\endlastfoot
\tablestripes
Extraction $\to$ Quality & 0.72** & 0.78** & 0.68** & 0.81** & 0.70** & 0.74** & 0.85** \\
Quality $\to$ Performance & 0.68** & 0.74** & 0.65** & 0.78** & 0.66** & 0.71** & 0.82** \\
Selection $\to$ Performance & 0.45** & 0.52** & 0.42** & 0.55** & 0.48** & 0.46** & 0.58** \\
Governance $\to$ Trustworthiness & 0.61** & 0.58** & 0.64** & 0.62** & 0.59** & 0.63** & 0.65** \\
Explainability $\to$ Adoption & 0.54** & 0.51** & 0.57** & 0.56** & 0.52** & 0.55** & 0.48** \\
\end{xltabular}}
\vspace{0.5em}\noindent\textit{Note.} ** $p<0.01$ (all paths significant). PLS-SEM with bootstrap resampling (1,000 resamples). Feature Quality $\rightarrow$ Model Performance is the strongest path for ASD ($\beta$ range: 0.65--0.82).

\subsection{R$^2$ and Effect Sizes}

Table~\ref{tab:phase5_r2_f2} documents the explained variance ($R^2$) for each endogenous construct and the Cohen $f^2$ effect sizes for each structural path across all seven disease domains.
\needspace{8\baselineskip}
{\tablestripes\tablefontsize
\begin{xltabular}{\textwidth}{@{} >{\raggedright\arraybackslash}p{3.5cm} R R R R R R R @{}}
\caption[SEM R-squared and Effect Sizes per Disease Domain]{SEM Explained Variance (R$^2$) and Effect Sizes ($f^2$) per Disease Domain}\label{tab:phase5_r2_f2} \\
\toprule
\thead{Construct} & \thead{SZ} & \thead{EP} & \thead{ST} & \thead{ASD} & \thead{PD} & \thead{DEP} & \thead{CA} \\
\midrule
\endfirsthead
\multicolumn{8}{@{}l}{\textit{(\,Tab.~\ref{tab:phase5_r2_f2} continued from previous page\,)}} \\
\toprule
\thead{Construct} & \thead{SZ} & \thead{EP} & \thead{ST} & \thead{ASD} & \thead{PD} & \thead{DEP} & \thead{CA} \\
\midrule
\endhead
\bottomrule
\endlastfoot
\tablestripes
R$^2$ (Feat.\ Quality) & 0.52 & 0.61 & 0.46 & 0.66 & 0.49 & 0.55 & 0.72 \\
R$^2$ (Model Perf.) & 0.68 & 0.75 & 0.62 & 0.80 & 0.65 & 0.70 & 0.84 \\
R$^2$ (Trustworthiness) & 0.37 & 0.34 & 0.41 & 0.38 & 0.35 & 0.40 & 0.42 \\
$f^2$ (Extr.\ $\to$ Quality) & 0.38 & 0.45 & 0.32 & 0.52 & 0.36 & 0.41 & 0.58 \\
$f^2$ (Quality $\to$ Perf.) & 0.42 & 0.51 & 0.35 & 0.58 & 0.39 & 0.45 & 0.62 \\
$f^2$ (Gov.\ $\to$ Trust) & 0.28 & 0.25 & 0.31 & 0.29 & 0.26 & 0.30 & 0.33 \\
\end{xltabular}}
\vspace{0.5em}\noindent\textit{Note.} $f^2$ thresholds: 0.02 small, 0.15 medium, 0.35 large \citep{cohen1988statistical}. All Feature Quality $\rightarrow$ Model Performance paths show large effect sizes ($f^2 > 0.35$), confirming that feature quality is the dominant predictor of classification performance.

\subsection{Hypothesis Validation Summary}

% MOVED TO CHAPTER 4 MAIN TEXT (Mar 2026)
\iffalse
\needspace{8\baselineskip}
\noindent Table~\ref{tab:app_phase5_hypothesis} summarises hypothesis validation results --- support status across all seven disease domains for appendix review.

{\tablestripes\tablefontsize
\begin{xltabular}{\textwidth}{@{} c L l l @{}}
\caption[Hypothesis Validation Results]{Hypothesis Validation Results --- Support Status Across All Seven Disease Domains}\label{tab:app_phase5_hypothesis} \\
\toprule
\thead{H} & \thead{Hypothesis Statement} & \thead{Result} & \thead{Domains Supported} \\
\midrule
\endfirsthead
\multicolumn{4}{@{}l}{\textit{(\,Tab.~\ref{tab:app_phase5_hypothesis} continued from previous page\,)}} \\
\toprule
\thead{H} & \thead{Hypothesis Statement} & \thead{Result} & \thead{Domains Supported} \\
\midrule
\endhead
\bottomrule
\endlastfoot
\tablestripes
H1 & Systematic feature extraction significantly improves feature quality & Supported & 7/7 ($p<0.01$) \\
H2 & Higher feature quality leads to superior model performance & Supported & 7/7 ($p<0.01$) \\
H3 & Multi-method feature selection enhances classification accuracy & Supported & 7/7 ($p<0.01$) \\
H4 & AI governance integration improves model trustworthiness & Supported & 7/7 ($p<0.01$) \\
H5 & Explainability mechanisms increase clinical adoption readiness & Supported & 7/7 ($p<0.01$) \\
\end{xltabular}}
\vspace{0.5em}\noindent\textit{Note.} All five hypotheses are supported across all seven disease domains, providing strong evidence for the RGAIG framework's generalisability.
\fi

% [SPACE-FIX] clearpage removed to eliminate empty page

%% ======================================================================
%% G.6  PHASE 6 — MEDIATION + MODERATION (Per Disease)
%% ======================================================================
\section{Phase 6 --- Mediation and Moderation Analysis}
\label{sec:phase6_medmod}

Phase 6 tests indirect effects (mediation: does feature quality mediate the extraction-performance relationship?) and conditional effects (moderation: does domain complexity moderate the quality-performance link?).

\subsection{Mediation Analysis Results}

Table~\ref{tab:phase6_mediation} presents the direct, indirect, and total effects for the mediation path (Feature Extraction $\rightarrow$ Feature Quality $\rightarrow$ Model Performance) across all seven disease domains, distinguishing partial from full mediation.
\needspace{8\baselineskip}
{\tablestripes\tablefontsize
\begin{xltabular}{\textwidth}{@{} l R R R l @{}}
\caption[Mediation Analysis Results per Disease Domain]{Mediation Analysis --- Direct, Indirect, and Total Effects per Disease Domain}\label{tab:phase6_mediation} \\
\toprule
\thead{Disease} & \thead{Direct Effect} & \thead{Indirect Effect} & \thead{Total Effect} & \thead{Mediation Type} \\
\midrule
\endfirsthead
\multicolumn{5}{@{}l}{\textit{(\,Tab.~\ref{tab:phase6_mediation} continued from previous page\,)}} \\
\toprule
\thead{Disease} & \thead{Direct Effect} & \thead{Indirect Effect} & \thead{Total Effect} & \thead{Mediation Type} \\
\midrule
\endhead
\bottomrule
\endlastfoot
\tablestripes
Schizophrenia & 0.28* & 0.42** & 0.70** & Partial \\
Epilepsy & 0.22* & 0.51** & 0.73** & Partial \\
Stress & 0.31* & 0.38** & 0.69** & Partial \\
ASD & 0.18 & 0.58** & 0.76** & Full \\
Parkinson & 0.25* & 0.44** & 0.69** & Partial \\
Depression & 0.27* & 0.45** & 0.72** & Partial \\
Cancer & 0.15 & 0.62** & 0.77** & Full \\
\end{xltabular}}
\vspace{0.5em}\noindent\textit{Note.} * $p<0.05$, ** $p<0.01$. Mediation path: Feature Extraction $\rightarrow$ Feature Quality $\rightarrow$ Model Performance. Bootstrap resampling (5,000 resamples). ASD and Cancer show full mediation (direct effect non-significant), indicating feature quality completely explains the extraction-performance link.

\subsection{Moderation Analysis Results}

Table~\ref{tab:phase6_moderation} reports the interaction effects of domain complexity as a moderator on the feature quality--performance link, including the incremental variance explained ($\Delta R^2$) per disease domain.
\needspace{8\baselineskip}
{\tablestripes\tablefontsize
\begin{xltabular}{\textwidth}{@{} l R R R l @{}}
\caption[Moderation Analysis Results per Disease Domain]{Moderation Analysis --- Interaction Effects of Domain Complexity on Feature Quality-Performance Link}\label{tab:phase6_moderation} \\
\toprule
\thead{Disease} & \thead{Main Effect ($\beta$)} & \thead{Interaction ($\beta$)} & \thead{$\Delta$R$^2$} & \thead{Significant?} \\
\midrule
\endfirsthead
\multicolumn{5}{@{}l}{\textit{(\,Tab.~\ref{tab:phase6_moderation} continued from previous page\,)}} \\
\toprule
\thead{Disease} & \thead{Main Effect ($\beta$)} & \thead{Interaction ($\beta$)} & \thead{$\Delta$R$^2$} & \thead{Significant?} \\
\midrule
\endhead
\bottomrule
\endlastfoot
\tablestripes
Schizophrenia & 0.68** & 0.18* & 0.03 & Yes \\
Epilepsy & 0.74** & 0.22** & 0.05 & Yes \\
Stress & 0.65** & 0.15* & 0.02 & Yes \\
ASD & 0.78** & 0.12 & 0.01 & No \\
Parkinson & 0.66** & 0.20* & 0.04 & Yes \\
Depression & 0.71** & 0.16* & 0.03 & Yes \\
Cancer & 0.82** & 0.08 & 0.01 & No \\
\end{xltabular}}
\vspace{0.5em}\noindent\textit{Note.} Moderator: domain complexity (number of channels $\times$ sampling rate $\times$ artefact prevalence). Interaction significant for 5/7 domains. ASD and Cancer show no significant moderation, suggesting their classification performance is robust regardless of domain complexity.

\subsection{Mediation-Moderation Model}
\needspace{20\baselineskip}
\begin{figure}[!htbp]
\centering
\resizebox{0.97\textwidth}{!}{%
\begin{tikzpicture}[
    construct/.style={rectangle, rounded corners=4pt, draw=ggunavy, fill=ggunavy!8, minimum width=2.8cm, minimum height=1cm, font=\small\bfseries, text=ggunavy, align=center},
    mod/.style={rectangle, rounded corners=4pt, draw=orange!60, fill=orange!10, minimum width=2.4cm, minimum height=1.0cm, font=\small\bfseries, text=orange!70!black, align=center},
    arr/.style={-{Stealth[length=4mm, width=3mm]}, very thick, ggunavy},
    modarr/.style={-{Stealth[length=4mm, width=3mm]}, very thick, orange!60, dashed}
]
\node[construct] (ext) at (0,0) {Feature\\Extraction};
\node[construct] (qual) at (5.5,0) {Feature\\Quality};
\node[construct] (perf) at (11,0) {Model\\Performance};
\node[mod] (comp) at (8.25,2.5) {Domain\\Complexity};

%% Direct and indirect paths
\draw[arr] (ext) -- node[above, font=\small] {$a$ = 0.68--0.85} (qual);
\draw[arr] (qual) -- node[above, font=\small] {$b$ = 0.65--0.82} (perf);
\draw[arr, ggunavy!40] (ext) to[bend right=25] node[below, font=\small] {$c'$ = 0.15--0.31} (perf);

%% Moderation
\draw[modarr] (comp) -- node[right, font=\small, text=orange!70!black] {$\beta$ = 0.08--0.22} (8.25,0.5);

%% Labels
\node[font=\small, text=ggunavy!60, below] at (2.75,-0.7) {Indirect: $a \times b$ = 0.38--0.62};
\node[font=\small, text=ggunavy!60, below] at (5.5,-1.2) {Mediation: Partial (5/7) or Full (2/7)};
\end{tikzpicture}%
}%
\caption[Mediation-moderation model]{Integrated mediation-moderation model. Feature quality mediates the extraction-performance relationship (indirect effects range 0.38--0.62). Domain complexity moderates the quality-performance path in 5/7 domains (dashed orange arrow). Path $c'$ represents the residual direct effect after mediation.}
\label{fig:phase6_medmod}
\end{figure}

% [SPACE-FIX] clearpage removed to eliminate empty page

%% ======================================================================
%% G.7  PHASE 7 — QCA + TISM (Per Disease)
%% ======================================================================
\section{Phase 7 --- Alternative Models (QCA + TISM)}
\label{sec:phase7_qca_tism}

Phase 7 validates the robustness of findings through alternative analytical approaches: fuzzy-set Qualitative Comparative Analysis (fsQCA) identifies multiple causal pathways, while Total Interpretive Structural Modelling (TISM) establishes hierarchical relationships.

\subsection{fsQCA Truth Table and Causal Paths}

Table~\ref{tab:phase7_qca} enumerates the five equifinal causal configurations identified through fuzzy-set Qualitative Comparative Analysis, together with their consistency, coverage, and the disease domains each path covers.
\needspace{8\baselineskip}
{\tablestripes\tablefontsize
\begin{xltabular}{\textwidth}{@{} c l R R L @{}}
\caption[fsQCA Causal Paths to High Classification Accuracy]{fsQCA Causal Paths --- Configurations Leading to High Classification Accuracy ($>$90\%)}\label{tab:phase7_qca} \\
\toprule
\thead{Path} & \thead{Configuration} & \thead{Consistency} & \thead{Coverage} & \thead{Domains} \\
\midrule
\endfirsthead
\multicolumn{5}{@{}l}{\textit{(\,Tab.~\ref{tab:phase7_qca} continued from previous page\,)}} \\
\toprule
\thead{Path} & \thead{Configuration} & \thead{Consistency} & \thead{Coverage} & \thead{Domains} \\
\midrule
\endhead
\bottomrule
\endlastfoot
\tablestripes
P1 & High Feature Quality + High Freq. Features & 0.92 & 0.68 & ASD, EP, SZ, DEP \\
P2 & High Feature Quality + Strong Governance & 0.89 & 0.55 & All 7 domains \\
P3 & High Connectivity + High Spatial Features & 0.85 & 0.42 & SZ, ASD, DEP \\
P4 & High Explainability + High Feature Selection & 0.87 & 0.48 & EP, PD, ST \\
P5 & High Radiomics + Strong Governance & 0.94 & 0.31 & Cancer only \\
\end{xltabular}}
\vspace{0.5em}\noindent\textit{Note.} Consistency threshold: $>$0.80. Coverage indicates proportion of outcome explained. Path P2 (Feature Quality + Governance) achieves the broadest coverage, supporting the RGAIG framework's dual emphasis on technical quality and governance.

\subsection{TISM Hierarchical Structure}

Table~\ref{tab:phase7_tism} maps the five-level TISM hierarchy, from foundation activities (data preparation and feature extraction) through to top-level outcomes (clinical adoption and model trustworthiness).
\needspace{8\baselineskip}
{\tablestripes\tablefontsize
\begin{xltabular}{\textwidth}{@{} c L L @{}}
\caption[TISM Level Partitioning Results]{TISM Level Partitioning --- Hierarchical Structure of Framework Components}\label{tab:phase7_tism} \\
\toprule
\thead{Level} & \thead{Variables} & \thead{Interpretation} \\
\midrule
\endfirsthead
\multicolumn{3}{@{}l}{\textit{(\,Tab.~\ref{tab:phase7_tism} continued from previous page\,)}} \\
\toprule
\thead{Level} & \thead{Variables} & \thead{Interpretation} \\
\midrule
\endhead
\bottomrule
\endlastfoot
\tablestripes
1 (Top/Outcome) & Clinical Adoption, Model Trustworthiness & Final outcomes driven by all lower levels \\
2 & Classification Accuracy, Explainability & Intermediate outcomes that enable adoption \\
3 & Feature Quality, Model Performance & Technical achievements that enable accuracy \\
4 & Feature Selection, Governance Integration & Methodological choices that ensure quality \\
5 (Foundation) & Data Preparation, Feature Extraction & Base-level activities that everything depends on \\
\end{xltabular}}
\vspace{0.5em}\noindent\textit{Note.} Level 5 (foundation) must be achieved before higher levels. This hierarchy validates the 10-phase framework's sequential design: data preparation and feature extraction precede all downstream analysis.

\subsection{TISM Hierarchical Model}
\needspace{20\baselineskip}
\begin{figure}[!htbp]
\centering
\resizebox{0.97\textwidth}{!}{%
\begin{tikzpicture}[
    level/.style={rectangle, rounded corners=4pt, draw=ggunavy, minimum width=10cm, minimum height=0.9cm, font=\small\bfseries, text=ggunavy, align=center},
    arr/.style={-{Stealth[length=4mm, width=3mm]}, very thick, ggunavy}
]
\node[level, fill=red!10] (l1) at (0,0) {Level 1: Clinical Adoption + Model Trustworthiness};
\node[level, fill=orange!10] (l2) at (0,-1.5) {Level 2: Classification Accuracy + Explainability};
\node[level, fill=yellow!12] (l3) at (0,-3) {Level 3: Feature Quality + Model Performance};
\node[level, fill=green!10] (l4) at (0,-4.5) {Level 4: Feature Selection + Governance Integration};
\node[level, fill=blue!10] (l5) at (0,-6) {Level 5: Data Preparation + Feature Extraction};

\draw[arr] (l5) -- (l4);
\draw[arr] (l4) -- (l3);
\draw[arr] (l3) -- (l2);
\draw[arr] (l2) -- (l1);

%% Interpretive labels
\node[font=\small, text=ggunavy!60, right] at (5.5,-0.75) {Enables clinical deployment};
\node[font=\small, text=ggunavy!60, right] at (5.5,-2.25) {Requires explainable results};
\node[font=\small, text=ggunavy!60, right] at (5.5,-3.75) {Depends on feature quality};
\node[font=\small, text=ggunavy!60, right] at (5.5,-5.25) {Built on clean data};
\end{tikzpicture}%
}%
\caption[TISM hierarchical model]{TISM hierarchical model showing five levels of the RGAIG framework. Foundation activities (data preparation, feature extraction) must be completed before feature quality can be assessed, which in turn enables classification accuracy and explainability, ultimately driving clinical adoption and model trustworthiness.}
\label{fig:phase7_tism}
\end{figure}

% [SPACE-FIX] clearpage removed to eliminate empty page

%% ======================================================================
%% G.8  PHASE 8 — QUALITATIVE ANALYSIS
%% ======================================================================
\section{Phase 8 --- Qualitative Analysis}
\label{sec:phase8_qualitative}

Phase 8 provides the interpretive ``why'' behind the quantitative findings, using thematic coding and discourse analysis applied to the biomarker patterns and clinical implications of each disease domain.

\subsection{Thematic Coding Results}

Table~\ref{tab:phase8_themes} summarises the primary qualitative theme and key clinical insight extracted from the biomarker analysis of each disease domain.
\needspace{8\baselineskip}
{\tablestripes\tablefontsize
\begin{xltabular}{\textwidth}{@{} l L L @{}}
\caption[Thematic Coding Results by Disease Domain]{Thematic Coding --- Key Themes Identified from Per-Disease Biomarker Analysis}\label{tab:phase8_themes} \\
\toprule
\thead{Disease} & \thead{Primary Theme} & \thead{Key Insight} \\
\midrule
\endfirsthead
\multicolumn{3}{@{}l}{\textit{(\,Tab.~\ref{tab:phase8_themes} continued from previous page\,)}} \\
\toprule
\thead{Disease} & \thead{Primary Theme} & \thead{Key Insight} \\
\midrule
\endhead
\bottomrule
\endlastfoot
\tablestripes
Schizophrenia & Alpha deficit with preserved connectivity & Alpha power reduction ($-$38\%) is the strongest discriminator; gamma coherence preserved in early stages \\
Epilepsy & Ictal delta spikes with temporal localisation & Delta power surges ($>$300\% increase) precisely localise seizure foci; line length correlates with seizure severity \\
Stress & Multimodal arousal signature & EDA (skin conductance) + reduced HRV together discriminate stress better than EEG alone; validates wearable integration \\
ASD & Frontal-posterior connectivity deficit & Reduced F-P coherence ($-$45\%) is the most consistent ASD marker; age-dependent progression observed \\
Parkinson & Beta suppression with tremor-band peaks & Resting-state beta suppression and characteristic 4--6 Hz tremor band peak co-occur; medication state modulates effect \\
Depression & Left-lateralised alpha asymmetry & Frontal alpha asymmetry index correlates with depression severity ($r = -0.62$); treatment response predictable \\
Cancer & Texture heterogeneity as malignancy marker & GLCM contrast and entropy discriminate benign from malignant; shape irregularity adds specificity \\
\end{xltabular}}

\subsection{Triangulation with Quantitative Findings}

% MOVED TO CHAPTER 4 MAIN TEXT (Mar 2026)
\iffalse
\needspace{8\baselineskip}
\noindent Table~\ref{tab:app_phase8_triangulation} summarises triangulation --- alignment of quantitative results with qualitative themes per disease do for appendix review.

{\tablestripes\tablefontsize
\begin{xltabular}{\textwidth}{@{} l L L l @{}}
\caption[Triangulation of Quantitative and Qualitative Findings]{Triangulation --- Alignment of Quantitative Results with Qualitative Themes per Disease Domain}\label{tab:app_phase8_triangulation} \\
\toprule
\thead{Disease} & \thead{Quantitative Finding} & \thead{Qualitative Insight} & \thead{Alignment} \\
\midrule
\endfirsthead
\multicolumn{4}{@{}l}{\textit{(\,Tab.~\ref{tab:app_phase8_triangulation} continued from previous page\,)}} \\
\toprule
\thead{Disease} & \thead{Quantitative Finding} & \thead{Qualitative Insight} & \thead{Alignment} \\
\midrule
\endhead
\bottomrule
\endlastfoot
\tablestripes
SZ & 97.17\% accuracy; alpha power top feature & Alpha deficit is primary biomarker & Strong \\
EP & 99.02\% accuracy; delta power dominant & Ictal delta spikes are pathognomonic & Strong \\
ST & 94.17\% accuracy; multimodal features & EDA+HRV combination is essential & Strong \\
ASD & 97.67\% accuracy; connectivity top-3 & F-P coherence deficit well-established & Strong \\
PD & 100\% accuracy; beta suppression key & Motor beta deficit is classical PD marker & Strong \\
DEP & 91.07\% accuracy; alpha asymmetry & Left-lateralised alpha deficit validated & Strong \\
CA & 99.02\% accuracy; texture features & GLCM contrast is gold-standard radiomics & Strong \\
\end{xltabular}}
\vspace{0.5em}\noindent\textit{Note.} All seven domains show strong alignment between quantitative model features and established clinical biomarker knowledge, validating the clinical interpretability of the RGAIG framework.
\fi

% [SPACE-FIX] clearpage removed to eliminate empty page

%% ======================================================================
%% G.9  PHASE 9 — INTEGRATION (Final Model)
%% ======================================================================
\section{Phase 9 --- Integration and Final Model}
\label{sec:phase9_integration}

Phase 9 synthesises all findings from Phases 1--8 into a unified framework, validates through triangulation, and establishes the final RGAIG model.

\subsection{Result Consolidation}

Table~\ref{tab:phase9_consolidation} consolidates the key finding from each of the nine preceding analysis phases into a single summary view.
\needspace{8\baselineskip}
{\tablestripes\tablefontsize
\begin{xltabular}{\textwidth}{@{} l l L @{}}
\caption[Consolidated Results Across All Analysis Phases]{Consolidated Results --- Key Findings from Each Analysis Phase}\label{tab:phase9_consolidation} \\
\toprule
\thead{Phase} & \thead{Source} & \thead{Key Finding} \\
\midrule
\endfirsthead
\multicolumn{3}{@{}l}{\textit{(\,Tab.~\ref{tab:phase9_consolidation} continued from previous page\,)}} \\
\toprule
\thead{Phase} & \thead{Source} & \thead{Key Finding} \\
\midrule
\endhead
\bottomrule
\endlastfoot
\tablestripes
1 & Data Prep & All 7 datasets cleaned; non-normal distributions confirmed \\
2 & Descriptive & Frequency-domain features show highest discriminative power \\
3 & Reliability & All constructs exceed CR $>$0.85, AVE $>$0.50 \\
4 & Factor Analysis & 3--5 factors per domain; CFA fit indices all acceptable \\
5 & SEM & All 5 hypotheses supported across all 7 domains \\
6 & Mediation & Feature quality mediates extraction$\rightarrow$performance (5 partial, 2 full) \\
7 (QCA) & Alt. Models & 5 causal paths identified; P2 covers all 7 domains \\
7 (TISM) & Alt. Models & 5-level hierarchy; data prep is foundation \\
8 & Qualitative & All domains show strong quant-qual alignment \\
\end{xltabular}}

\subsection{Multi-Method Triangulation}

Table~\ref{tab:phase9_triangulation} cross-validates the five core findings against four analytical methods (SEM, QCA, TISM, qualitative), demonstrating convergent support across all approaches.
\needspace{8\baselineskip}
{\tablestripes\tablefontsize
\begin{xltabular}{\textwidth}{@{} L l l l l @{}}
\caption[Multi-Method Triangulation Validation]{Multi-Method Triangulation --- Cross-Validation of Findings Across Analytical Approaches}\label{tab:phase9_triangulation} \\
\toprule
\thead{Finding} & \thead{SEM} & \thead{QCA} & \thead{TISM} & \thead{Qual.} \\
\midrule
\endfirsthead
\multicolumn{5}{@{}l}{\textit{(\,Tab.~\ref{tab:phase9_triangulation} continued from previous page\,)}} \\
\toprule
\thead{Finding} & \thead{SEM} & \thead{QCA} & \thead{TISM} & \thead{Qual.} \\
\midrule
\endhead
\bottomrule
\endlastfoot
\tablestripes
Feature quality drives performance & Supported & In all paths & Level 3 & Confirmed \\
Governance enhances trustworthiness & Supported & Path P2 & Level 4 & Confirmed \\
Explainability enables adoption & Supported & Path P4 & Level 2 & Confirmed \\
Data preparation is foundational & Baseline & Assumed & Level 5 & Confirmed \\
Multiple paths to high accuracy & --- & 5 paths & --- & Confirmed \\
\end{xltabular}}
\vspace{0.5em}\noindent\textit{Note.} All five core findings are validated by at least three analytical methods, demonstrating strong methodological triangulation and robustness.

\subsection{Theoretical and Practical Contributions}

Table~\ref{tab:phase9_contributions} enumerates the eight theoretical and practical contributions (C1--C8) that emerge from the integrated analysis of all preceding phases.
\needspace{8\baselineskip}
{\tablestripes\tablefontsize
\begin{xltabular}{\textwidth}{@{} c l L @{}}
\caption[Theoretical and Practical Contributions]{Theoretical and Practical Contributions of the RGAIG Framework}\label{tab:phase9_contributions} \\
\toprule
\thead{ID} & \thead{Type} & \thead{Contribution} \\
\midrule
\endfirsthead
\multicolumn{3}{@{}l}{\textit{(\,Tab.~\ref{tab:phase9_contributions} continued from previous page\,)}} \\
\toprule
\thead{ID} & \thead{Type} & \thead{Contribution} \\
\midrule
\endhead
\bottomrule
\endlastfoot
\tablestripes
C1 & Theoretical & Unified ASD-focused AI framework validated across 7 disease domains \\
C2 & Theoretical & Empirical validation of feature quality as mediator (not just predictor) \\
C3 & Theoretical & TISM hierarchical model establishing data preparation as foundation \\
C4 & Theoretical & Multi-path causal model (QCA) showing equifinal routes to accuracy \\
C5 & Practical & Reusable feature engineering pipeline for new disease domains \\
C6 & Practical & Governance integration framework for clinical AI deployment \\
C7 & Practical & Explainability mechanisms that increase clinical adoption readiness \\
C8 & Practical & 525-module quality taxonomy for AI system validation \\
\end{xltabular}}

\subsection{Final Integrated Framework}
\needspace{20\baselineskip}
\begin{figure}[!htbp]
\centering
\resizebox{0.97\textwidth}{!}{%
\begin{tikzpicture}[
    layer/.style={rectangle, rounded corners=4pt, draw=ggunavy, minimum width=12cm, minimum height=1.2cm, font=\small\bfseries, text=ggunavy, align=center},
    arr/.style={-{Stealth[length=4mm, width=3mm]}, very thick, ggunavy},
    feedback/.style={-{Stealth[length=4mm, width=3mm]}, very thick, ggunavy!40, dashed}
]
%% Layers
\node[layer, fill=blue!10] (input) at (0,-6) {Input Layer: Data Preparation + Feature Extraction + Feature Selection};
\node[layer, fill=green!10] (mech) at (0,-4) {Mechanism Layer: Feature Quality + Model Training + Validation};
\node[layer, fill=yellow!12] (output) at (0,-2) {Output Layer: Classification Accuracy + Explainability};
\node[layer, fill=orange!10] (govern) at (0,0) {Governance Layer: Responsible AI + Compliance + Trust};
\node[layer, fill=red!10] (adopt) at (0,2) {Adoption Layer: Clinical Deployment + Monitoring + Impact};

%% Arrows
\draw[arr] (input) -- (mech);
\draw[arr] (mech) -- (output);
\draw[arr] (output) -- (govern);
\draw[arr] (govern) -- (adopt);

%% Feedback
\draw[feedback] (adopt.east) -- ++(1.5,0) |- node[right, font=\small, pos=0.25, text=ggunavy!60] {Continuous improvement} (input.east);

%% Moderator
\node[rectangle, rounded corners=3pt, draw=orange!60, fill=orange!8, minimum width=3cm, font=\small\bfseries, text=orange!70!black, align=center] (mod) at (-7,-3) {Moderator:\\Domain\\Complexity};
\draw[-{Stealth[length=4mm,width=3mm]}, thick, orange!50, dashed, thick] (mod) -- (-3.5,-3);

%% QCA paths
\node[rectangle, rounded corners=3pt, draw=teal!60, fill=teal!8, minimum width=3cm, font=\small, text=teal!70!black, align=center] (qca) at (7,-3) {QCA: 5 equifinal\\paths to accuracy};
\draw[-{Stealth[length=4mm,width=3mm]}, thick, teal!50, dashed, thick] (qca) -- (3.5,-3);

\end{tikzpicture}%
}
\caption[Final integrated RGAIG framework]{Final integrated RGAIG framework showing five layers: input (data and features), mechanism (quality and training), output (accuracy and explainability), governance (RAI and compliance), and adoption (clinical deployment). Domain complexity moderates the mechanism-output link. QCA identifies five equifinal paths. A continuous improvement feedback loop connects adoption outcomes back to input refinement.}
\label{fig:phase9_framework}
\end{figure}

% [SPACE-FIX] clearpage removed to eliminate empty page

%% ======================================================================
%% G.10  PHASE 10 — CONCLUSION, LIMITATIONS & RECOMMENDATIONS
%% ======================================================================
\section{Phase 10 --- Conclusion, Limitations, and Recommendations}
\label{sec:phase10_conclusion}

Phase 10 synthesises the complete research journey, validates research objectives, acknowledges limitations, and provides actionable recommendations.

\subsection{Research Objective Validation}

Table~\ref{tab:phase10_objectives} maps each of the four research objectives (O1--O4) to its achievement status and the supporting evidence drawn from the analysis phases.
\needspace{8\baselineskip}
{\tablestripes\tablefontsize
\begin{xltabular}{\textwidth}{@{} c L l L @{}}
\caption[Research Objective Validation Status]{Research Objective Validation --- Achievement Status and Supporting Evidence}\label{tab:phase10_objectives} \\
\toprule
\thead{Obj.} & \thead{Objective Statement} & \thead{Status} & \thead{Evidence} \\
\midrule
\endfirsthead
\multicolumn{4}{@{}l}{\textit{(\,Tab.~\ref{tab:phase10_objectives} continued from previous page\,)}} \\
\toprule
\thead{Obj.} & \thead{Objective Statement} & \thead{Status} & \thead{Evidence} \\
\midrule
\endhead
\bottomrule
\endlastfoot
\tablestripes
O1 & Develop a unified AI framework for ASD diagnosis & Achieved & RGAIG framework validated across 7 domains \\
O2 & Validate feature engineering pipeline across heterogeneous datasets & Achieved & 42 features/domain avg.; all domains $>$91\% accuracy \\
O3 & Integrate AI governance and responsible AI principles & Achieved & 525-module quality taxonomy; 10 governance pillars \\
O4 & Demonstrate clinical adoption readiness through explainability & Achieved & SHAP/LIME integration; SEM path $\beta>$0.48 in all domains \\
\end{xltabular}}

\subsection{Limitations}

Table~\ref{tab:phase10_limitations} documents the six research limitations (L1--L6), their scope boundaries, and the mitigation strategies or justifications applied to each.
\needspace{8\baselineskip}
{\tablestripes\tablefontsize
\begin{xltabular}{\textwidth}{@{} c L L @{}}
\caption[Research Limitations and Mitigation]{Research Limitations --- Scope Boundaries and Mitigation Strategies}\label{tab:phase10_limitations} \\
\toprule
\thead{\#} & \thead{Limitation} & \thead{Mitigation / Justification} \\
\midrule
\endfirsthead
\multicolumn{3}{@{}l}{\textit{(\,Tab.~\ref{tab:phase10_limitations} continued from previous page\,)}} \\
\toprule
\thead{\#} & \thead{Limitation} & \thead{Mitigation / Justification} \\
\midrule
\endhead
\bottomrule
\endlastfoot
\tablestripes
L1 & Sample sizes vary across domains (48--19,847) & Bootstrap resampling (1,000 resamples) ensures robust inference even for small samples; still exceeds minimum for SEM \\
L2 & Cross-sectional design limits temporal inference & Longitudinal validation identified as future work; current cross-validation provides within-study robustness \\
L3 & Public datasets may not reflect clinical populations & Datasets sourced from established repositories (PhysioNet, TCIA, OpenNeuro) with documented protocols; generalisability acknowledged as limitation \\
L4 & Single framework applied to heterogeneous modalities & Modality-specific branches (EEG vs imaging) address heterogeneity while maintaining unified governance layer \\
L5 & No prospective clinical trial validation & Retrospective analysis with k-fold/LOSO CV serves as proxy; clinical trial recommended as future work \\
L6 & Explainability assessed technically, not with end-users & User-facing explainability study recommended as future work; technical SHAP/LIME validates model interpretability \\
\end{xltabular}}

\subsection{Future Research Recommendations}

Table~\ref{tab:phase10_future} presents seven future research directions (F1--F7) for extending and validating the RGAIG framework beyond the scope of this dissertation.
\needspace{8\baselineskip}
{\tablestripes\tablefontsize
\begin{xltabular}{\textwidth}{@{} c L L @{}}
\caption[Future Research Recommendations]{Future Research Recommendations --- Directions for Extending the RGAIG Framework}\label{tab:phase10_future} \\
\toprule
\thead{\#} & \thead{Recommendation} & \thead{Rationale} \\
\midrule
\endfirsthead
\multicolumn{3}{@{}l}{\textit{(\,Tab.~\ref{tab:phase10_future} continued from previous page\,)}} \\
\toprule
\thead{\#} & \thead{Recommendation} & \thead{Rationale} \\
\midrule
\endhead
\bottomrule
\endlastfoot
\tablestripes
F1 & Longitudinal validation with temporal data & Track model drift and stability over time; validate PSI/CSI monitoring \\
F2 & Multi-centre clinical trial & Validate generalisability across institutions and populations \\
F3 & Additional disease domains (Alzheimer's, ADHD) & Extend framework to neurodegenerative and neurodevelopmental conditions \\
F4 & Federated learning integration & Enable multi-site model training without data sharing; address privacy \\
F5 & Real-time edge deployment & Validate on Emotiv EPOC X and wearable devices in clinical settings \\
F6 & End-user explainability study & Assess clinician comprehension and trust with human factors evaluation \\
F7 & Regulatory pathway mapping & Map framework outputs to FDA/CE-MDR/HIPAA requirements \\
\end{xltabular}}

\subsection{Consolidated Statement}

The RGAIG framework demonstrates that a unified, governance-integrated AI architecture can achieve clinically meaningful diagnostic accuracy across seven heterogeneous disease domains. Feature quality, mediated through systematic extraction and selection, is the primary driver of classification performance. Governance and explainability are not afterthoughts but structural enablers of clinical adoption. The framework's generalisability is validated through multi-method triangulation (SEM, QCA, TISM, and qualitative analysis), with all five hypotheses supported across all seven domains.

\textbf{The central thesis argument is confirmed: responsible generative AI governance, operationalised through the RGAIG framework, transforms ASD AI from a technical achievement into a clinically deployable, trustworthy, and governable system.}

\subsection{Phase 10 Quality Assessment}

Table~\ref{tab:phase10_quality_appendix} provides the final quality assessment checklist, confirming that the conclusion phase satisfies all eight evaluation criteria.
\needspace{8\baselineskip}
{\tablestripes\tablefontsize
\begin{xltabular}{\textwidth}{@{} l L l @{}}
\caption[Phase 10 Quality Assessment]{Phase 10 Quality Assessment --- Final Validation Checklist}\label{tab:phase10_quality_appendix} \\
\toprule
\thead{Area} & \thead{Assessment Criterion} & \thead{Status} \\
\midrule
\endfirsthead
\multicolumn{3}{@{}l}{\textit{(\,Tab.~\ref{tab:phase10_quality_appendix} continued from previous page\,)}} \\
\toprule
\thead{Area} & \thead{Assessment Criterion} & \thead{Status} \\
\midrule
\endhead
\bottomrule
\endlastfoot
\tablestripes
Alignment & Conclusions match stated research objectives (O1--O4) & Pass \\
Clarity & Findings expressed in accessible language & Pass \\
Contribution & Academic and practical contributions clearly stated (C1--C8) & Pass \\
Honesty & Limitations acknowledged with mitigation strategies (L1--L6) & Pass \\
Impact & Practical recommendations actionable for stakeholders & Pass \\
Completeness & All 7 disease domains covered in conclusion & Pass \\
Triangulation & Multi-method validation documented and aligned & Pass \\
Future work & Specific, feasible recommendations provided (F1--F7) & Pass \\
\end{xltabular}}

\clearpage
